\chapter{Linear Approximation and Differentials}

In the previous chapter, we calculated tangent slopes for implicit curves
and related rates. But why spend so much effort finding slopes? Because a
tangent line is the simplest, most useful linear model of the curve itself.

There is one idea hiding underneath every derivative in this book, and this
chapter is where we finally say it out loud:

\begin{center}
\emph{A smooth curve, looked at closely enough, is a straight line.}
\end{center}

That insight forms the foundation of differential calculus. Geometrically, the derivative
$f'(a)$ is the slope of the
\emph{tangent line that the curve approaches} under magnification near $x=a$. Once
you recognize that smooth means locally linear, the surrounding ideas --- tangent
lines, error estimates, Newton's method, and the differentials $\dd y$ and $\dd x$
used throughout physics --- become straightforward applications of replacing a curve
by its tangent line and bounding the remainder.

This chapter makes that replacement precise. We will build the
\emph{linearization} of a function, measure the error you commit by using it,
pin down what $\dd y$ and $\dd x$ \emph{really mean} (they are not zero,
and they are not ``infinitely small''), use them to track how measurement errors
propagate through a formula, and take a first look at Newton's method, which is
just linear approximation applied over and over.

\chapterline

\section{Zooming In: Why Smooth Means Straight}
\label{ch14:zooming}

Take any function you like that is differentiable at a point -- say
$f(x)=x^2$ near $x=1$ -- and imagine doing the following experiment. Plot it.
Then zoom in on the point $(1,1)$: cut the viewing window in half, then in half
again, then again. What happens to the parabola?

It flattens. The more you zoom, the straighter the visible piece of curve looks,
until on a small enough window you genuinely cannot tell it from a straight line
drawn with a ruler. Figure~\ref{ch14:fig-zoom} shows three frames of this
experiment.

\begin{figure}[H]
\centering
\begin{tikzpicture}
\begin{axis}[
  width=4.6cm, height=4.6cm,
  title={\footnotesize window $\pm 1$},
  title style={font=\small},
  xmin=0, xmax=2, ymin=-0.2, ymax=4,
  xtick={0,1,2}, ytick={0,2,4},
  tick label style={font=\tiny},
  axis lines=left,
]
\addplot[curveblue, thick, samples=80, domain=0:2] {x^2};
\addplot[tangentred, thick, domain=0:2] {2*x-1};
\addplot[mark=*, mark size=1pt, black] coordinates {(1,1)};
\end{axis}
\end{tikzpicture}
\hspace{0.2cm}
\begin{tikzpicture}
\begin{axis}[
  width=4.6cm, height=4.6cm,
  title={\footnotesize window $\pm 0.3$},
  title style={font=\small},
  xmin=0.7, xmax=1.3, ymin=0.4, ymax=1.7,
  xtick={0.7,1,1.3}, ytick={0.5,1,1.5},
  tick label style={font=\tiny},
  axis lines=left,
]
\addplot[curveblue, thick, samples=80, domain=0.7:1.3] {x^2};
\addplot[tangentred, thick, domain=0.7:1.3] {2*x-1};
\addplot[mark=*, mark size=1pt, black] coordinates {(1,1)};
\end{axis}
\end{tikzpicture}
\hspace{0.2cm}
\begin{tikzpicture}
\begin{axis}[
  width=4.6cm, height=4.6cm,
  title={\footnotesize window $\pm 0.07$},
  title style={font=\small},
  xmin=0.93, xmax=1.07, ymin=0.85, ymax=1.16,
  xtick={0.93,1,1.07}, ytick={0.9,1,1.15},
  tick label style={font=\tiny},
  axis lines=left,
]
\addplot[curveblue, thick, samples=80, domain=0.93:1.07] {x^2};
\addplot[tangentred, thick, domain=0.93:1.07] {2*x-1};
\addplot[mark=*, mark size=1pt, black] coordinates {(1,1)};
\end{axis}
\end{tikzpicture}
\caption{Zooming in on $y=x^2$ at $(1,1)$. The curve (blue) and its tangent line
$y=2x-1$ (red) become indistinguishable as the window shrinks. ``Smooth''
literally means ``looks straight up close.''}
\label{ch14:fig-zoom}
\end{figure}

\begin{intuitionbox}
Differentiability \emph{is} this zoom property. To say $f$ has a derivative at
$a$ is to say: there is a single straight line through $(a,f(a))$ that the curve
approaches under magnification. A function with a corner (like $\abs{x}$ at $0$) fails
the test -- no matter how far you zoom into the corner, you still see two
different slopes meeting, never one line. Smoothness is the condition that the zoom
converges to a single line.
\end{intuitionbox}

Why should you care that the curve becomes straight? Because straight lines are
the only functions we can compute with easily. A line is just ``start
here, go up by the slope times how far you walk.'' No squaring, no exponentials,
no trig tables. If we are willing to accept a small error, we can replace any
hard smooth function, \emph{near a point}, by a line and do arithmetic instead of
analysis. That trade -- a little error in exchange for enormous simplicity -- is
the engine of applied mathematics.

\begin{walkbox}
Picture the surface of the Earth. Locally -- in your backyard, on a soccer field
-- it looks flat, and flat-earth geometry (Pythagoras, straight survey lines)
works beautifully. The error only shows up over hundreds of miles, where the
curvature accumulates. The Earth is the function; the flat backyard is its
tangent plane; the error is the curvature. Where else in life do you replace a
curved thing by a flat one because the flat version is good enough up close?
\end{walkbox}

\section{The Tangent Line as the Best Linear Approximation}
\label{ch14:tangent-best}

Let us turn the picture into algebra. We have a function $f$ differentiable at a
point $a$. We want the straight line that best matches $f$ near $a$. A straight
line through the point $(a,f(a))$ has the form
\[
L(x) = f(a) + m\,(x-a),
\]
and the only freedom we have is the slope $m$. Which slope makes $L$ hug the
curve best?

Geometrically the answer is obvious from the zoom picture: pick the slope of the
curve itself at $a$, namely $m=f'(a)$. That gives the \emph{tangent line}. But
``looks best'' is not a definition. Let us say precisely what ``best'' means and
then prove the tangent line wins.

\begin{definition}[Linearization]
\label{ch14:def-linearization}
Let $f$ be differentiable at $a$. The \textbf{linearization} of $f$ at $a$ is the
linear function
\[
L(x) \;=\; f(a) + f'(a)\,(x-a).
\]
The approximation $f(x)\approx L(x)$ for $x$ near $a$ is called the
\textbf{linear} (or \textbf{tangent-line}) \textbf{approximation}.
\end{definition}

Notice the anatomy of $L$: it carries the function's \emph{height} $f(a)$ and its
\emph{slope} $f'(a)$ at the single point $a$, and nothing else. It is the most a
straight line can know about $f$. Read it aloud: ``Start at the known value
$f(a)$. Walk a distance $(x-a)$. For each unit you walk, climb by the rate
$f'(a)$.''

\subsection{In what sense is it ``best''?}

Compare the tangent line to any other line through $(a,f(a))$ with a different
slope $m\neq f'(a)$. The error of a line $L_m(x)=f(a)+m(x-a)$ is
\[
E_m(x) = f(x) - L_m(x).
\]
At $x=a$ every such line is exact: $E_m(a)=0$, because they all pass through the
point. So the contest is about how fast the error grows as you step away. Write
$x = a+h$ and divide the error by the step $h$:
\[
\frac{E_m(a+h)}{h} = \frac{f(a+h)-f(a)}{h} - m.
\]
As $h\to 0$ the difference quotient $\frac{f(a+h)-f(a)}{h}$ tends to $f'(a)$.
Therefore
\[
\frac{E_m(a+h)}{h} \;\longrightarrow\; f'(a)-m.
\]
Look at what that limit says. If $m\neq f'(a)$, it is a nonzero number; the
error is roughly $(f'(a)-m)\cdot h$, growing \emph{in proportion to} $h$. But if
$m=f'(a)$, the limit is $0$: the error, divided by $h$, approaches zero. The error of the
tangent line shrinks \emph{faster than $h$ itself}. That is the precise sense in
which the tangent line is the uniquely best linear approximation.

\begin{focusbox}
\textbf{The defining property of the tangent line.} Among all lines through
$(a,f(a))$, the tangent line is the only one whose error is \emph{small compared
to the step}: writing $x=a+h$,
\[
\lim_{h\to 0}\frac{f(a+h) - \big[f(a)+f'(a)\,h\big]}{h} \;=\; 0.
\]
Every other line leaves an error proportional to $h$; the tangent line leaves an
error that is negligible next to $h$. We say the tangent line agrees with $f$ ``to
first order.''
\end{focusbox}

\begin{whybox}
Why does matching the slope kill the first-order error? Because the derivative
\emph{is} the leading-order rate of change. Near $a$,
\[
f(a+h) = f(a) + f'(a)\,h + (\text{error}),
\]
and the very definition of $f'(a)$ as a limit forces the error to be small
relative to $h$. Choosing the slope to be anything other than $f'(a)$ leaves a
leftover term $(f'(a)-m)h$ that is itself proportional to $h$ -- a first-order
mistake that swamps the curvature. Matching the slope cancels the leading first-order discrepancy, leaving an error
that scales with $h^2$ as curvature pulls the curve away from the line.
\end{whybox}

\subsection{A first worked example}

\begin{example}[Approximating a square root]
\label{ch14:ex-sqrt}
Estimate $\sqrt{4.1}$ without a calculator.

Let $f(x)=\sqrt{x}$, and choose the \emph{anchor} $a=4$ because we know
$\sqrt{4}=2$ exactly and $4$ is close to $4.1$. We need the slope:
\[
f'(x) = \frac{1}{2\sqrt{x}}, \qquad f'(4) = \frac{1}{2\cdot 2} = \frac14.
\]
The linearization is
\[
L(x) = f(4) + f'(4)(x-4) = 2 + \tfrac14 (x-4).
\]
Plug in $x=4.1$:
\[
\sqrt{4.1} \approx L(4.1) = 2 + \tfrac14 (0.1) = 2 + 0.025 = 2.025.
\]
The true value is $\sqrt{4.1} = 2.024845\ldots$, so our estimate is right to four
decimal places. We did it with one division.
\end{example}

Picking a good anchor $a$ means balancing two things: you want a point where $f$ and $f'$
can be computed exactly, and which sits as close to your target as possible.
Here $a=4$ serves both purposes: $4$ is a perfect square and lies adjacent to $4.1$.

\begin{labbox}
Estimate $\sqrt{4.4}$ the same way: $L(4.4) = 2 + \tfrac14(0.4) = 2.1$. The true
value is $2.0976\ldots$. Notice the error has grown from about $0.00015$ (at the
step $h=0.1$) to about $0.0024$ (at the step $h=0.4$): the step grew by a factor
of $4$, and the error grew by a factor of about $16=4^2$. That quadratic growth
reflects the curvature, and Section~\ref{ch14:error} explains
why the error scales with $h^2$.
\end{labbox}

\begin{recoverybox}
\textbf{Pause and think.} Look at our estimate $\sqrt{4.1} \approx 2.025$. Before
consulting a calculator: is $2.025$ slightly too large or slightly too small?
Look at the graph of $y=\sqrt{x}$ in Figure~\ref{ch14:fig-sqrt}. As $x$ grows,
the curve gets flatter---it bends downward (concave down). Any tangent line
drawn to a downward-bending curve must ride \emph{above} the curve. Therefore,
without calculating a single extra digit, you can know with total certainty that
$2.025$ is an overestimate. Check: $2.025^2 = 4.100625 > 4.1$. The geometry
never lies.
\end{recoverybox}

\begin{figure}[H]
\centering
\begin{tikzpicture}
\begin{axis}[
  width=10cm, height=6cm,
  xlabel={$x$}, ylabel={$y$},
  xmin=2.5, xmax=6, ymin=1.4, ymax=2.6,
  xtick={3,4,5,6}, ytick={1.5,2,2.5},
  axis lines=left,
  legend pos=south east, legend style={font=\footnotesize},
  samples=100,
]
\addplot[curveblue, thick, domain=2.5:6] {sqrt(x)};
\addlegendentry{$y=\sqrt{x}$}
\addplot[tangentred, thick, domain=2.5:6] {2 + 0.25*(x-4)};
\addlegendentry{$L(x)=2+0.25(x-4)$}
\addplot[mark=*, black, mark size=1.3pt] coordinates {(4,2)};
\addplot[dashed, gray] coordinates {(4.1,1.4) (4.1,2.025)};
\addplot[mark=*, tangentred, mark size=1.2pt] coordinates {(4.1,2.025)};
\node[font=\footnotesize, anchor=north west] at (axis cs:4.15,1.95) {$L(4.1)=2.025$};
\node[font=\footnotesize, anchor=south east] at (axis cs:3.95,2.05) {anchor at $a=4$};
\end{axis}
\end{tikzpicture}
\caption{The tangent line to $\sqrt{x}$ at $a=4$. Near the anchor the red line
and blue curve are nearly identical; the gap (the error) opens slowly as you move
away, and because $\sqrt{x}$ bends downward, the line sits slightly above the
curve.}
\label{ch14:fig-sqrt}
\end{figure}

\begin{trapbox}
Linearization is a \emph{local} tool. It is only trustworthy near the anchor. If
you used the line $L(x)=2+\tfrac14(x-4)$ to estimate $\sqrt{100}$ you would get
$2+\tfrac14(96)=26$, which is absurd (the truth is $10$). The approximation does
not ``know'' anything about the curve far from $a$ -- it only ever knew the height
and slope at one point. Always ask: is my target genuinely close to my anchor?
\end{trapbox}

\section{What the Differentials \texorpdfstring{$\dd y$}{dy} and \texorpdfstring{$\dd x$}{dx} Really Mean}
\label{ch14:differentials}

You have been writing $\dv{y}{x}$ for chapters now, defined rigorously as the limit of a difference
quotient. That care was necessary while constructing the foundations. Now we can clarify the
underlying geometric picture: there are two distinct quantities, $\dd x$ and $\dd y$, whose ratio
equals the derivative when defined on the tangent line. Leibniz originally thought of $\dd x$ as
an infinitesimal change. While intuitive, that notion required two centuries of refinement to place
on solid ground. Here is the modern formulation.

\begin{definition}[The differentials $\dd x$ and $\dd y$]
\label{ch14:def-differentials}
Let $y=f(x)$ be differentiable at $x$. We declare $\dd x$ to be an
\textbf{independent real variable} -- an ordinary number, which you are free to
choose, representing a step in the $x$-direction. Then we \emph{define} the
\textbf{differential of $y$} by
\[
\dd y \;=\; f'(x)\,\dd x.
\]
$\dd y$ is the rise of the \emph{tangent line} over a run of $\dd x$.
\end{definition}

Read that definition slowly, because it dissolves a confusion you may have
carried for years.

\begin{happeningbox}
\textbf{What $\dd x$ and $\dd y$ really are.} They are not zero. They are not
``infinitely small.'' $\dd x$ is just a number you pick -- the size of a step.
And $\dd y$ is the corresponding \emph{rise along the tangent line}, not along the
curve. Look at Figure~\ref{ch14:fig-dydx}. Take a real step $\dd x = \Delta x$
from the point $(a,f(a))$. Two things rise:
\begin{itemize}
\item the \emph{curve} rises by $\Delta y = f(a+\Delta x) - f(a)$ -- the true
change;
\item the \emph{tangent line} rises by $\dd y = f'(a)\,\dd x$ -- the linear
estimate of the change.
\end{itemize}
The differential $\dd y$ is the \textbf{approximation}; the increment $\Delta y$
is the \textbf{truth}. The gap between them, $\Delta y - \dd y$, is the error of
linear approximation -- exactly the quantity we studied in the last section. And
now the notation $\dv{y}{x}=f'(x)$ is literally a quotient of two numbers $\dd y$
and $\dd x$, because that is how we defined $\dd y$. Leibniz's fraction was real
all along; we just had to say which line the rise is measured on.
\end{happeningbox}

\begin{figure}[H]
\centering
\begin{tikzpicture}[scale=1.05]
% curve: f(x)=0.35 x^2 + 0.4 ; anchor a=1.4
\draw[curveblue, thick, samples=100, domain=0.4:3.0, variable=\x]
  plot (\x, {0.35*\x*\x + 0.4});
\def\a{1.4}      % anchor x
\def\fa{1.086}   % f(a) = 0.35*1.96+0.4
\def\fp{0.98}    % f'(a) = 0.7*1.4
\def\xb{2.7}     % a + dx
% tangent line at a
\draw[tangentred, thick] (0.6,{\fa+\fp*(0.6-\a)}) -- (3.0,{\fa+\fp*(3.0-\a)});
\filldraw[black] (\a,\fa) circle (1.6pt);
\node[below left, font=\small] at (\a,\fa) {$(a,f(a))$};
% curve and tangent points at x=xb
% f(xb)=0.35*7.29+0.4=2.9515 ; tangent: fa+fp*(xb-a)=1.086+0.98*1.3=2.360
\filldraw[curveblue] (\xb,2.9515) circle (1.4pt);
\filldraw[tangentred] (\xb,2.360) circle (1.4pt);
% horizontal dx
\draw[densely dashed, gray!70] (\a,\fa) -- (\xb+0.85,\fa);
\draw[densely dashed, gray!70] (\xb,2.9515) -- (\xb+0.85,2.9515);
\draw[<->] (\a,0.2) -- (\xb,0.2) node[midway, below=2pt, font=\small] {$\dd x=\Delta x$};
\draw[densely dashed] (\xb,2.9515) -- (\xb,0.2);
% dy (tangent rise)
\draw[<->, tangentred] (\xb+0.08,\fa) -- (\xb+0.08,2.360);
\node[right, tangentred, font=\small] at (\xb+0.08,{(\fa+2.360)/2}) {$\dd y$};
% Delta y (true rise)
\draw[<->, curveblue] (\xb+0.7,\fa) -- (\xb+0.7,2.9515);
\node[right, curveblue, font=\small] at (\xb+0.7,{(\fa+2.9515)/2}) {$\Delta y$};
% error gap
\draw[<->, structurepurple, thick] (\xb-0.22,2.360) -- (\xb-0.22,2.9515);
\node[left, structurepurple, font=\footnotesize] at (\xb-0.22,{(2.360+2.9515)/2}) {error};
\node[font=\footnotesize, curveblue] at (0.95,2.7) {$y=f(x)$};
\node[font=\footnotesize, tangentred] at (3.0,3.2) {tangent};
\end{tikzpicture}
\caption{The true meaning of differentials. Over a real step $\dd x=\Delta x$,
the tangent line rises by $\dd y=f'(a)\,\dd x$ while the curve rises by the true
amount $\Delta y$. The differential $\dd y$ estimates $\Delta y$; their
difference (purple) is the approximation error, governed by the curvature.}
\label{ch14:fig-dydx}
\end{figure}

\begin{historybox}
Calculus began not with limits, but with the practical struggle to approximate curved motion using straight lines. In the 1660s and 1670s, Isaac Newton and Gottfried Wilhelm Leibniz independently discovered that the slope of a curve could be found by imagining a change so small that the curve was effectively straight. Newton framed this dynamically, speaking of ``fluxions'' (instantaneous velocities) and ``evanescent increments'' (quantities caught at the fleeting instant of vanishing). Leibniz developed the enduring notation of differentials: he pictured $\dd x$ as an \emph{infinitesimal}---a quantity smaller than any positive real number you can name, yet strictly greater than zero. With $\dd y = f'(x)\,\dd x$, Leibniz treated derivatives as genuine ratios of geometric lengths, which made the chain rule $\dv{y}{x}=\dv{y}{u}\dv{u}{x}$ look like simple fraction cancellation.

Their practical results were accurate, but their foundations provoked immediate debate. In 1734, Bishop George Berkeley published his critique \emph{The Analyst}, targeting a clear logical inconsistency: in calculating a derivative, mathematicians treated $\dd x$ as nonzero so they could divide by it, and then set $\dd x = 0$ to discard the remaining terms. Within the ordinary real numbers, no single quantity can satisfy both conditions.

The response took over a century to mature. In 1768, Leonhard Euler published his \emph{Institutiones calculi differentialis}, demonstrating that stepping along tangent lines ($y_{n+1} = y_n + f'(x_n)\Delta x$) provided a universal numerical tool---now called Euler's method---for solving otherwise intractable differential equations.

The full theoretical resolution arrived in 1823 with Augustin-Louis Cauchy. Cauchy rebuilt the derivative as the limit of a difference quotient of finite real numbers. Cauchy also showed that linear approximation carries a rigorous quantitative guarantee: using the Mean Value Theorem, he proved that the error $R_1(x) = f(x) - L(x)$ is bounded by the maximum curvature: $|R_1(x)| \le \frac{1}{2} M (x-a)^2$, where $|f''(t)| \le M$. Leibniz's $\dd x$ and $\dd y$ were restored as ordinary, finite numbers: $\dd x$ is a chosen horizontal step, and $\dd y = f'(a)\,\dd x$ is the exact vertical rise of the tangent line.

The story held one final surprise. In the 1960s, the logician Abraham Robinson used mathematical logic to prove that Leibniz's dream of true, consistent infinitesimals was logically sound after all, founding \emph{nonstandard analysis}. Today, while the theory of single-variable linear approximation is complete, a vital question remains open in applied mathematics: determining exactly when and where the linearization of massive, complex systems---such as climate dynamics, turbulent fluids, or deep neural networks---breaks down.
\end{historybox}

\subsection{Differentials as a notation for approximation}

The single most useful sentence about differentials is this:
\[
\boxed{\;\Delta y \;\approx\; \dd y \;=\; f'(x)\,\dd x\;}
\]
``The true change is approximately the linear change.'' Everything in the rest of
this chapter is an application of that boxed line. Let us redo
Example~\ref{ch14:ex-sqrt} in this language to see how compact it becomes.

\begin{example}[$\sqrt{4.1}$ via differentials]
\label{ch14:ex-sqrt-diff}
We want $\sqrt{4.1}$. Anchor at $x=4$, take the step $\dd x = 0.1$. With
$y=\sqrt{x}$ we have $\dd y = \frac{1}{2\sqrt{x}}\,\dd x$, so
\[
\dd y = \frac{1}{2\sqrt 4}(0.1) = \frac{0.1}{4} = 0.025.
\]
Then
\[
\sqrt{4.1} = \sqrt 4 + \Delta y \approx \sqrt 4 + \dd y = 2 + 0.025 = 2.025,
\]
matching what we found before. The differential $\dd y=0.025$ \emph{is} the
estimated change in $\sqrt{x}$ caused by nudging $x$ from $4$ to $4.1$.
\end{example}

\begin{connectionbox}
The boxed approximation $\Delta y \approx f'(x)\,\dd x$ is the one-variable seed
of the \emph{total differential} you will meet in multivariable calculus:
\[
\dd z = \pdv{f}{x}\,\dd x + \pdv{f}{y}\,\dd y.
\]
There, a small move in the plane changes $z$ by (rate in $x$)$\times$(step in
$x$) plus (rate in $y$)$\times$(step in $y$): the tangent \emph{plane} replaces
the tangent line, but the idea -- locally, smooth equals linear -- is
identical. Get it down in one variable and the multivariable version costs you
almost nothing.
\end{connectionbox}

\section{Linearization and Error: How Wrong Are You?}
\label{ch14:error}

An approximation you cannot bound is a guess. The reason linear approximation is
trustworthy -- rather than merely convenient -- is that we can say precisely how
large the error is. Define the error of the linearization at $a$ as
\[
E(x) = f(x) - L(x) = f(x) - \big[f(a)+f'(a)(x-a)\big].
\]
We already know (Section~\ref{ch14:tangent-best}) that $E$ vanishes faster than
the step: $E(a)=0$ and $E(x)/(x-a)\to 0$. That tells us the error is small, but
not \emph{how} small. The next theorem gives the size exactly:
\emph{the error is controlled by the second derivative}---by the
curvature that the straight-line approximation omits.

\begin{theorem}[Error of linear approximation]
\label{ch14:thm-error}
Suppose $f$ has a continuous second derivative on an interval containing $a$ and
$x$. Then there is a point $c$ between $a$ and $x$ such that
\[
f(x) - \big[f(a) + f'(a)(x-a)\big] \;=\; \frac{f''(c)}{2}\,(x-a)^2.
\]
\end{theorem}

Before the proof, absorb the shape of the result. Writing $h = x-a$ for the step:
\[
E = \tfrac12\, f''(c)\, h^2.
\]
Three things to notice. First, the error is proportional to $h^2$, not $h$ -- so
halving your distance from the anchor quarters the error, exactly the square-law
we saw numerically in the lab box. Second, the multiplier is $f''$, the
curvature: a flatter curve (small $f''$) is approximated better, and a perfectly
straight $f$ (with $f''=0$) has \emph{zero} error, as it must. Third, the sign of
$f''$ tells you which way you missed: if the curve is concave up ($f''>0$) the
tangent line lies below it and you underestimate; concave down ($f''<0$) and you
overestimate. That is why in Figure~\ref{ch14:fig-sqrt} the line for the
downward-bending $\sqrt{x}$ sat \emph{above} the curve, and our estimate $2.025$
came out slightly high. Figure~\ref{ch14:fig-errgrow} shows this $h^2$ growth
directly: the true error tracks a parabola bounded by the guaranteed bound.

\begin{whybox}
Where does the $\tfrac12$ come from, and why $f''$? Think of building $f$ back up
from its derivatives, like reconstructing a journey from your speedometer. The
value $f(a)$ is where you start. The term $f'(a)h$ is how far you would travel if
your speed never changed -- the linear estimate. But your speed \emph{does} change,
at the rate $f''$ (that is what the second derivative is: the rate the slope
changes). Over a step $h$, a steadily changing speed contributes an
\emph{average} extra displacement. Just as distance under constant acceleration is
$\tfrac12 a t^2$ in physics, the curvature contributes $\tfrac12 f'' h^2$. The
$\tfrac12$ is the same $\tfrac12$ that appears in $\tfrac12 at^2$ -- it is the
average of a quantity ramping linearly from $0$ to its full value. The error term
is the first correction you would add to make the straight-line estimate curved.
\end{whybox}

\begin{proof}[Proof of Theorem~\ref{ch14:thm-error}]
This proof relies on a fundamental result that we will systematically develop and
prove in the very next chapter: the \emph{Mean Value Theorem}, which guarantees
that a smooth function's instantaneous rate of change matches its average rate of
change across an interval at some intermediate point. Here we use both the
standard Mean Value Theorem and its two-function extension, the Cauchy Mean Value
Theorem (which asserts that the ratio of overall changes $\frac{g(x)-g(a)}{p(x)-p(a)}$
equals the ratio of instantaneous derivatives $\frac{g'(c_1)}{p'(c_1)}$ at some point $c_1$).
You can follow the core idea without waiting for Chapter~15: we compare the growth
of the error to the growth of the quadratic step $(x-a)^2$.

Fix $a$ and let $x$ vary; write $h=x-a$ and consider the two functions
\[
g(t) = f(t) - \big[f(a)+f'(a)(t-a)\big],
\qquad
p(t) = (t-a)^2,
\]
both on the interval between $a$ and $x$. We want $g(x)/p(x)$. Notice
$g(a)=0$ and $p(a)=0$, and also
\[
g'(t) = f'(t) - f'(a), \qquad p'(t) = 2(t-a),
\]
so $g'(a)=0$ and $p'(a)=0$ as well. By the Cauchy mean value theorem applied to
$g$ and $p$ on the interval from $a$ to $x$ (using $g(a)=p(a)=0$), there is a
point $c_1$ between $a$ and $x$ with
\[
\frac{g(x)}{p(x)} = \frac{g(x)-g(a)}{p(x)-p(a)} = \frac{g'(c_1)}{p'(c_1)}
= \frac{f'(c_1)-f'(a)}{2(c_1-a)}.
\]
Now apply the ordinary mean value theorem to $f'$ on the interval from $a$ to
$c_1$: there is a point $c$ between $a$ and $c_1$ (hence between $a$ and $x$) with
\[
\frac{f'(c_1)-f'(a)}{c_1-a} = f''(c).
\]
Substituting,
\[
\frac{g(x)}{p(x)} = \frac{f''(c)}{2}
\quad\Longrightarrow\quad
g(x) = \frac{f''(c)}{2}\,(x-a)^2,
\]
which is exactly the claim, since $g(x)=E(x)$.
\end{proof}

\begin{focusbox}
\textbf{Usable error bound.} You rarely know the mystery point $c$, but you do not
need it. If $\abs{f''(t)}\le M$ for all $t$ between $a$ and $x$, then
\[
\big|f(x)-L(x)\big| \;\le\; \frac{M}{2}\,(x-a)^2.
\]
Replace the unknown $f''(c)$ by an \emph{upper bound} on the curvature and you get
a guaranteed, conservative error bar.
\end{focusbox}

\begin{example}[Putting a number on the $\sqrt{4.1}$ error]
\label{ch14:ex-sqrt-bound}
For $f(x)=\sqrt{x}$ we have $f''(x) = -\tfrac14 x^{-3/2}$. On the interval
$[4,4.1]$ the largest size of $f''$ occurs at the left end $x=4$ (where
$x^{-3/2}$ is biggest):
\[
\abs{f''(x)} \le \tfrac14 \cdot 4^{-3/2} = \tfrac14\cdot \tfrac18 = \tfrac{1}{32}
= M.
\]
With $h = 0.1$,
\[
\big|\sqrt{4.1}-2.025\big| \le \frac{M}{2}h^2
= \frac{1}{64}\,(0.1)^2 = \frac{0.01}{64} \approx 0.000156.
\]
The actual error was $0.000154\ldots$ -- our bound is essentially tight, and we
obtained it knowing only that $\sqrt{x}$ does not bend too sharply on $[4,4.1]$.
Because $f''<0$ there, we even know the \emph{sign}: the estimate is high, as it
was.
\end{example}

\begin{figure}[H]
\centering
\begin{tikzpicture}
\begin{axis}[
  width=10cm, height=5.6cm,
  xlabel={step $h=x-a$}, ylabel={error $\big|f(x)-L(x)\big|$},
  xmin=0, xmax=1, ymin=0, ymax=0.18,
  xtick={0,0.25,0.5,0.75,1}, ytick={0,0.05,0.1,0.15},
  axis lines=left,
  legend pos=north west, legend style={font=\footnotesize},
  samples=120,
]
% true error of sqrt(4+h) linearized at a=4: sqrt(4+h) - (2 + h/4)
\addplot[curveblue, thick, domain=0:1] {2 + 0.25*x - sqrt(4+x)};
\addlegendentry{true error of $\sqrt{x}$ at $a=4$}
% quadratic envelope M/2 h^2 with M=1/32
\addplot[tangentred, thick, dashed, domain=0:1] {(1/64)*x^2};
\addlegendentry{$\tfrac{M}{2}h^2=\tfrac{1}{64}h^2$}
\end{axis}
\end{tikzpicture}
\caption{The linearization error grows like $h^2$. The blue curve is the actual
error of approximating $\sqrt{x}$ by its tangent line at $a=4$; the dashed red
parabola is the guaranteed bound $\tfrac{M}{2}h^2$ from
Theorem~\ref{ch14:thm-error}. The error hugs a parabola, not a line -- doubling
the step roughly quadruples the error.}
\label{ch14:fig-errgrow}
\end{figure}

\begin{aestheticbox}
Theorem~\ref{ch14:thm-error} is the first step toward \emph{Taylor's theorem}, the
general approximation framework developed later. The pattern continues:
the linear approximation has error $\sim f'' h^2/2$; if you also match
the curvature you get the quadratic approximation, with error $\sim f''' h^3/6$;
and matching higher derivatives introduces higher powers of $h$. Linear
approximation is Taylor's theorem truncated at degree one.
\end{aestheticbox}

\section{Propagation of Error: Differentials in Practice}
\label{ch14:propagation}

Differentials provide a direct computational tool for laboratory sciences.
You measure something -- a radius, a voltage, a
time -- and your measurement has an uncertainty. You then \emph{compute} something
from it -- an area, a power, a distance. How does the uncertainty in what you
measured affect what you computed? The answer is one application of
$\dd y = f'(x)\,\dd x$, read in a new light: $\dd x$ is now the \emph{measurement
error}, and $\dd y$ is the resulting error in the computed quantity.

\begin{example}[Error in the area of a circle]
\label{ch14:ex-circle}
You measure the radius of a circular plate as $r = 10\,$cm, with your ruler good
to about $\dd r = \pm 0.05\,$cm. You report the area $A=\pi r^2$. What is the
uncertainty in $A$?

The area's differential is
\[
\dd A = \dv{A}{r}\,\dd r = 2\pi r\,\dd r.
\]
At $r=10$, $\dd r = 0.05$:
\[
\dd A = 2\pi (10)(0.05) = \pi \approx 3.14\ \text{cm}^2.
\]
So $A = \pi(10)^2 \pm \pi \approx 314.16 \pm 3.14\ \text{cm}^2$. A half-millimeter
slip in the radius costs you about $3\,$cm$^2$ of area.
\end{example}

\begin{figure}[H]
\centering
\begin{tikzpicture}[scale=0.95]
% Left: Circle with expanding ring
\begin{scope}[shift={(0,0)}]
  \fill[sourceblue!15] (0,0) circle (1.6);
  \fill[quotientorange!25, even odd rule] (0,0) circle (1.95) (0,0) circle (1.6);
  \draw[curveblue, thick] (0,0) circle (1.6);
  \draw[quotientorange, thick] (0,0) circle (1.95);
  \filldraw[black] (0,0) circle (1.4pt);
  
  \draw[->, thick, black] (0,0) -- (135:1.6) node[midway, above right, font=\footnotesize] {$r$};
  \draw[<->, thick, quotientorange!80!black] (1.6,0) -- (1.95,0) node[midway, below=2pt, font=\footnotesize] {$\dd r$};
  
  \node[curveblue, font=\footnotesize] at (0,-0.6) {$A = \pi r^2$};
  \node[quotientorange!80!black, font=\footnotesize, align=center] at (0,-2.5) {Added ring\\ true change $\Delta A$};
\end{scope}

% Arrow in between
\draw[->, very thick, gray!70] (2.4,0) -- (3.6,0) node[midway, above, font=\footnotesize] {unroll};

% Right: Unrolled strip
\begin{scope}[shift={(4.2,-0.8)}]
  \fill[quotientorange!25] (0,0) rectangle (5.2,0.7);
  \draw[quotientorange, thick] (0,0) rectangle (5.2,0.7);
  
  \fill[trapred!25] (5.2,0) -- (5.7,0.7) -- (5.2,0.7) -- cycle;
  \draw[trapred, thick] (5.2,0) -- (5.7,0.7) -- (5.2,0.7);
  
  \draw[<->, thick] (0,-0.25) -- (5.2,-0.25) node[midway, below, font=\footnotesize] {circumference $= 2\pi r$};
  \draw[<->, thick, quotientorange!80!black] (-0.25,0) -- (-0.25,0.7) node[midway, left, font=\footnotesize] {$\dd r$};
  
  \node[quotientorange!80!black, font=\footnotesize] at (2.6,0.35) {$\dd A = (2\pi r)\,\dd r$};
  \node[trapred, font=\scriptsize, align=left, anchor=west] at (5.8,0.35) {error $\pi(\dd r)^2$\\ (second order)};
  \node[font=\footnotesize, align=center] at (2.6,-1.7) {Linear estimate of added area:\\ height $\dd r$ times base $2\pi r$};
\end{scope}
\end{tikzpicture}
\caption{The geometry of error propagation for $A = \pi r^2$. A radius uncertainty $\dd r$
adds an annular ring of true area $\Delta A = \pi(r+\dd r)^2 - \pi r^2$. Unrolling the ring
yields a rectangle of length $2\pi r$ and width $\dd r$, giving the differential $\dd A =
2\pi r\,\dd r$. The tiny remaining red wedge has area $\pi(\dd r)^2$, which shrinks quadratically
and is ignored by the linear approximation.}
\label{ch14:fig-circle-error}
\end{figure}

\begin{recoverybox}
\textbf{Pause and think.} Suppose you measure the electrical resistance $R$ of a resistor with a $2\%$ uncertainty, and calculate the power dissipated $P = I^2 R$ at a known, fixed current $I$. What is the relative uncertainty in $P$? 
Now suppose instead that $R$ is known perfectly, but the current $I$ has a $2\%$ uncertainty. What is the uncertainty in $P = I^2 R$? 
Why does the current error hurt twice as much ($4\%$) as the resistance error?
\end{recoverybox}

Often the \emph{relative} error -- the fractional, or percentage, error -- is more
illuminating than the absolute error, because it is dimensionless and tells you
how many significant figures you can trust. Divide:
\[
\frac{\dd A}{A} = \frac{2\pi r\,\dd r}{\pi r^2} = 2\,\frac{\dd r}{r}.
\]

\begin{focusbox}
\textbf{Relative error doubles for area, triples for volume.} For $A = \pi r^2$,
\[
\frac{\dd A}{A} = 2\,\frac{\dd r}{r}.
\]
A $1\%$ error in radius becomes a $2\%$ error in area. For a volume $V = \tfrac43
\pi r^3$ the same calculation gives $\dfrac{\dd V}{V} = 3\,\dfrac{\dd r}{r}$ -- a
$1\%$ radius error becomes a $3\%$ volume error. The exponent in the power law
\emph{is} the error-amplification factor.
\end{focusbox}

\begin{whybox}
Why does the exponent become the multiplier? Because $\dd$(power) is just the
power rule turned into error language. If $y = x^n$ then $\dd y = n x^{n-1}\dd x$,
and dividing by $y=x^n$ gives
\[
\frac{\dd y}{y} = n\,\frac{\dd x}{x}.
\]
The logarithm makes it even cleaner: $\ln y = n\ln x$, so $\frac{\dd y}{y} =
n\,\frac{\dd x}{x}$ is exactly the differential of that identity (since
$\dd(\ln u) = \dd u/u$). Relative error is the \emph{differential of the
logarithm}. That is why, for any product or power expression, the relative errors
simply add up weighted by the exponents -- taking logs turns products into sums,
and differentiation turns the sums into added relative errors.
\end{whybox}

\begin{example}[Combining several measurements]
\label{ch14:ex-combine}
A pendulum's period is $T = 2\pi\sqrt{L/g}$. Suppose you know $L$ to within
$0.5\%$ and $g$ to within $0.2\%$. How accurately do you know $T$?

Take logarithms first: $\ln T = \ln(2\pi) + \tfrac12\ln L - \tfrac12 \ln g$.
Now take the differential of each side, using $\dd(\ln u)=\dd u/u$:
\[
\frac{\dd T}{T} = \tfrac12\,\frac{\dd L}{L} - \tfrac12\,\frac{\dd g}{g}.
\]
For a \emph{worst-case} bound we add the magnitudes:
\[
\left|\frac{\dd T}{T}\right| \le \tfrac12(0.5\%) + \tfrac12(0.2\%)
= 0.25\% + 0.1\% = 0.35\%.
\]
The logarithmic trick converted a messy square-root-and-quotient formula into a
one-line sum of weighted relative errors. This is the everyday error-propagation
arithmetic of physics labs.
\end{example}

\begin{applicationbox}
Error propagation via differentials runs through all of quantitative science.
A surveyor turning a measured angle into a distance, an engineer turning a
measured current into dissipated power $P=I^2R$ (so $\dd P/P = 2\,\dd I/I$ -- watch
your ammeter!), an economist turning a growth rate into a doubling time, a machine
learning engineer asking how a tiny perturbation $\dd x$ to an input shifts a
network's output -- that last one, the \emph{sensitivity} $\dd y = f'(x)\,\dd x$,
is precisely a single step of backpropagation. In every case the question ``how
much does the output wobble when the input wobbles?'' is answered by the
derivative times the input variation.
\end{applicationbox}

\section{The Newton Step: Linear Approximation in Reverse}
\label{ch14:newton}

So far we used the tangent line to estimate a function's \emph{value}. Now turn
the question inside out. Suppose we want to \emph{solve} an equation $f(x)=0$ --
to find a root -- but the equation is too hard to solve exactly. Examples are
everywhere: $\cos x = x$, or $x^5 - x - 1 = 0$, equations with no closed-form
solution. The tangent line provides a direct method:

\begin{center}
\emph{We cannot solve $f(x)=0$, but we can always solve $L(x)=0$,\\ because $L$
is a straight line.}
\end{center}

Start with a guess $x_0$ near the root. Replace the curve by its tangent line at
$x_0$ -- the best line we know how to build. Then find where \emph{that line}
crosses zero, and call it $x_1$. Because the line approximates the curve, its crossing
is generally much closer to the true root than $x_0$ was. Repeating this process
produces successive approximations that converge rapidly to the root.
Figure~\ref{ch14:fig-newton} shows the first two steps of this iteration.

\begin{figure}[H]
\centering
\begin{tikzpicture}
\begin{axis}[
  width=11cm, height=6.2cm,
  xlabel={$x$}, ylabel={$y$},
  xmin=0.6, xmax=3.2, ymin=-1.6, ymax=4.2,
  axis lines=middle,
  xtick={1,2,3}, ytick={0},
  samples=100,
  clip=false,
  legend pos=north west, legend style={font=\footnotesize},
]
% f(x)=x^2-2, root sqrt2 ~ 1.4142
\addplot[curveblue, thick, domain=0.6:3.2] {x^2-2};
\addlegendentry{$f(x)=x^2-2$}
% tangent at x0=2.5: zero at 1.65
\addplot[tangentred, thick, domain=1.65:2.5] {(2.5^2-2) + 2*2.5*(x-2.5)};
% tangent at x1=1.65: zero at ~1.431
\addplot[orange, thick, domain=1.431:1.65] {(1.65^2-2) + 2*1.65*(x-1.65)};
\addplot[mark=*, black, mark size=1.2pt, only marks] coordinates {(2.5,0) (1.65,0) (1.431,0)};
\node[below, font=\small] at (axis cs:2.5,-0.08) {$x_0$};
\node[below right, font=\small] at (axis cs:1.65,-0.05) {$x_1$};
\node[below left, font=\small] at (axis cs:1.431,-0.05) {$x_2$};
\addplot[densely dashed, gray] coordinates {(2.5,0) (2.5,4.25)};
\addplot[densely dashed, gray] coordinates {(1.65,0) (1.65,0.7225)};
\node[font=\footnotesize] at (axis cs:2.1,-1.0) {root $=\sqrt 2\approx 1.414$};
\end{axis}
\end{tikzpicture}
\caption{Newton's method on $f(x)=x^2-2$. From $x_0$, slide down the tangent
(red) to its zero $x_1$; from $x_1$, slide down the next tangent (orange) to
$x_2$. Each step lands far closer to the true root $\sqrt 2$.}
\label{ch14:fig-newton}
\end{figure}

Let us derive the step. The tangent line at $x_0$ is
\[
L(x) = f(x_0) + f'(x_0)(x - x_0).
\]
Set $L(x_1)=0$ and solve for $x_1$:
\[
0 = f(x_0) + f'(x_0)(x_1 - x_0)
\;\;\Longrightarrow\;\;
x_1 = x_0 - \frac{f(x_0)}{f'(x_0)}.
\]
This single formula drives the iteration.

\begin{focusbox}
\textbf{Newton's iteration.}
\[
x_{n+1} = x_n - \frac{f(x_n)}{f'(x_n)}.
\]
In words: from your current guess, step by ``minus the value over the slope.''
The value $f(x_n)$ tells you how far off you are; dividing by the slope
$f'(x_n)$ converts that vertical miss into the horizontal correction that a
straight line of that slope would need.
\end{focusbox}

\begin{example}[Computing $\sqrt 2$ by hand]
\label{ch14:ex-newton-sqrt2}
To find $\sqrt 2$, solve $f(x)=x^2-2=0$, with $f'(x)=2x$. The iteration becomes
\[
x_{n+1} = x_n - \frac{x_n^2 - 2}{2x_n} = \frac{x_n}{2} + \frac{1}{x_n},
\]
the famous ``average your guess with $2/\text{guess}$'' rule (known to the
Babylonians millennia before Newton). Start at $x_0 = 1.5$:
\[
\begin{aligned}
x_1 &= \tfrac{1.5}{2} + \tfrac{1}{1.5} = 0.75 + 0.66667 = 1.41667,\\
x_2 &= \tfrac{1.41667}{2} + \tfrac{1}{1.41667} \approx 0.70834 + 0.70588 = 1.41422,\\
x_3 &= 1.414214\ldots
\end{aligned}
\]
The true value is $\sqrt 2 = 1.4142136\ldots$. After just three steps we are
correct to seven digits. Notice the number of correct digits roughly
\emph{doubled} each step: $1$, then $3$, then $6$+. That explosive accuracy --
\emph{quadratic convergence} -- is the reward for using the best linear
approximation.
\end{example}

\begin{whybox}
Why does Newton's method converge so rapidly? Because the error it leaves
behind is exactly the linear-approximation error of Theorem~\ref{ch14:thm-error},
which is of order (current error)$^2$. Each step uses the tangent line, whose miss
is governed by $\tfrac12 f'' h^2$; so if your current distance to the root is $h$,
your next distance is of size $h^2$. Square a small number and it becomes tiny:
$0.01 \to 0.0001 \to 0.00000001$. The doubling of correct digits per step reflects this
$h \mapsto h^2$ scaling. Newton's method is fast for precisely the reason
the tangent line is accurate: matching the slope cancels the first-order error,
leaving only the squared remainder.
\end{whybox}

\begin{trapbox}
Newton's method is powerful but not foolproof. If $f'(x_n)=0$ the step divides by
zero -- the tangent is horizontal and never reaches the axis. A bad starting guess
can send the iterates flying away from the root, or trap them in a loop. And if
the root is one where $f'$ also vanishes (a double root), the fast doubling slows
to a crawl. We treat all of this carefully in the dedicated Newton's-method
chapter; here the point is only the \emph{idea} -- root-finding is linear
approximation run backwards, solving the line instead of the curve.
\end{trapbox}

\begin{intuitionbox}
\textbf{From algebra to geometry: Newton and Raphson.}
Newton described an early version of this method around 1669 for polynomial equations,
manipulating terms algebraically without explicitly picturing tangent slopes. In 1690,
Joseph Raphson streamlined the algebra into the iterative algorithm we use today, which
is why it is often called the \emph{Newton--Raphson method}. The modern formulation
$x_{n+1}=x_n-f(x_n)/f'(x_n)$, framed directly in terms of the derivative and sliding
down the tangent line, is the geometric insight that emerged once calculus matured.
It is remarkable that the same tangent line that approximates a curve's height also
hands us the premier equation-solver in scientific computing.
\end{intuitionbox}

\section{The Seed of All of Calculus}
\label{ch14:seed}

Look at what one idea -- \emph{locally, smooth is linear} -- has
already done in this single chapter.

\begin{itemize}
\item It \emph{defined} the derivative geometrically: $f'(a)$ is the slope of the
line the curve becomes under magnification.
\item It gave us \emph{linearization}: replace $f$ by $f(a)+f'(a)(x-a)$ to compute
hard values with easy arithmetic.
\item It pinned down \emph{differentials}: $\dd y = f'(x)\,\dd x$ is the tangent
line's rise, and $\dd y/\dd x$ is a true ratio.
\item It powered \emph{error propagation}: the same $\dd y = f'(x)\,\dd x$, read
as ``output wobble equals slope times input wobble.''
\item It produced \emph{Newton's method}: solve the tangent line instead of the
curve, and iterate.
\end{itemize}

And this is only the beginning of the book. The \emph{integral} is the same idea
inverted: chop a curved region into slivers so thin that each is a straight-sided
rectangle (smooth $=$ linear on a sliver), then add the slivers up. The
\emph{Fundamental Theorem} ties the slope picture to the area picture. A
\emph{differential equation} is a recipe that gives you the local slope $\dd
y/\dd x$ at every point and asks you to reassemble the global curve -- which you do
by stepping along tangent lines (that is literally Euler's method). Taylor series
push linear approximation to its limit, matching infinitely many derivatives.

\begin{aestheticbox}
One idea runs underneath all of it: \textbf{up close, a smooth function behaves like a straight line.} Differentiation is the act of \emph{finding} that local line; integration is the act of \emph{adding up} local linear pieces into a global total. Almost every tool in calculus -- the chain rule, substitution, Euler's method, Taylor series -- is another way of turning a curved problem into a locally straight one.
\end{aestheticbox}

\subsection{The Power, the Boundaries, and the Bridge to Chapter~15}
\label{ch14:what-unlocked}

\begin{connectionbox}
Linear approximation and differentials provide a universal trade: a tiny, controllable error in exchange for enormous computational simplicity.
\begin{itemize}
  \item \emph{Instant estimation}: Complicated values like $\sqrt{4.1}$ or $\ln(1.05)$ are evaluated with ordinary arithmetic: $f(a) + f'(a)\Delta x$.
  \item \emph{Sensitivity analysis}: The relation $\dd y = f'(a)\,\dd x$ tells scientists and engineers immediately how a measurement wobble $\dd x$ ripples into a computed result $\dd y$.
  \item \emph{Equation solving}: When an equation $f(x)=0$ defies an exact formula, we solve its straight-line stand-in $L(x)=0$ instead, yielding Newton's method.
\end{itemize}

Yet every straight line carries an inescapable blind spot: it has zero curvature. A tangent line assumes that the slope at the anchor point $a$ persists unchanged forever. The moment you step away from $a$, the curvature of the real curve pulls it away from the line, opening a quadratic error gap:
\[
  E(x) = f(x) - L(x) = \frac{f''(c)}{2}\,(x-a)^2.
\]
If a curve bends sharply (large $|f''|$) or if the step $h = x - a$ is large, the straight line quickly becomes inaccurate. A tangent line cannot turn corners, cannot track waves, and cannot anticipate asymptotes.

\textbf{Why we need the Mean Value Theorem next.}
Notice how heavily we leaned on Theorem~\ref{ch14:thm-error} to put guaranteed error bars on our estimates and to explain the quadratic convergence of Newton's method. But pause and ask: \emph{why are we guaranteed that such an intermediate point $c$ exists at all?} How can we be certain that a smooth curve cannot somehow wiggle wildly away from its tangent line between $a$ and $x$ without its second derivative registering that movement?

To turn linear error bounds from an intuitive rule of thumb into rigorous truth, we need a link connecting the local, instantaneous values of a derivative across an interval to the global change of the function over that interval. That link is the \textbf{Mean Value Theorem}, the central theorem of calculus that we explore in Chapter~15. Every rigorous error guarantee we used in this chapter is, in truth, a direct consequence of that one foundational theorem.

\textbf{Looking ahead: allowing approximations to bend.}
Once the Mean Value Theorem places our error bounds on firm ground, we will eventually generalize the tangent line. In Chapters~30 and~31 (\emph{Power Series and Taylor's Theorem}), we will allow our approximations to bend by matching the second derivative (curvature), third derivative, and beyond:
\[
  P_n(x) = f(a) + f'(a)(x-a) + \frac{f''(a)}{2!}(x-a)^2 + \cdots + \frac{f^{(n)}(a)}{n!}(x-a)^n.
\]
By incorporating higher powers of $(x-a)$, Taylor polynomials cling to curves across broad intervals. Linear approximation is the first, essential step on that journey: it is Taylor's theorem stopped at degree one.
\end{connectionbox}

\section{Exercises}
\label{ch14:exercises}

\begin{exercise}[Mental check: linearization at a glance]
Without picking up a pen or using paper, consider $f(x) = x^4$ anchored at $a = 1$.
\begin{enumerate}
\item What are the value $f(1)$ and the derivative $f'(1)$? Write the linearization $L(x)$ purely in your head.
\item Mentally estimate $(1.002)^4$.
\item Because $f''(x) = 12x^2 > 0$ for all $x$, does your mental estimate sit slightly above or slightly below the exact value?
\end{enumerate}
\end{exercise}

\begin{exercise}[Mental check: relative error scaling]
Answer each of the following using only the rule that the exponent is the error-amplification factor:
\begin{enumerate}
\item If you measure the side $x$ of a square with a $1.5\%$ uncertainty, what is the resulting percentage uncertainty in the computed area $A = x^2$?
\item If you measure the volume $V$ of a sphere and calculate its radius $r \propto V^{1/3}$, does the percentage error in $r$ grow or shrink compared to the error in $V$? By what factor?
\end{enumerate}
\end{exercise}

\begin{exercise}[Warm-up linearizations]
For each function, find the linearization $L(x)$ at the given anchor $a$, and use
it to estimate the requested value. Compare with a calculator.
\begin{enumerate}
\item $f(x)=\sqrt[3]{x}$ at $a=27$; estimate $\sqrt[3]{27.5}$.
\item $f(x)=\sin x$ at $a=0$; estimate $\sin(0.1)$ (radians).
\item $f(x)=\ln x$ at $a=1$; estimate $\ln(1.05)$.
\item $f(x)=(1+x)^{10}$ at $a=0$; estimate $(1.02)^{10}$.
\end{enumerate}
\end{exercise}

\begin{exercise}[Differentials directly]
A cube has side length measured as $s = 5.00\,$cm with possible error $\dd s =
\pm 0.02\,$cm. Use the differential of $V=s^3$ to estimate the largest possible
error in the computed volume, and the corresponding \emph{relative} error
$\dd V/V$. State the relative error as a percentage and confirm it equals
$3\,(\dd s/s)$.
\end{exercise}

\begin{exercise}[Sign of the error]
Without computing the true value, decide whether the linear approximation of
$f(x)=e^x$ at $a=0$ \emph{over}- or \emph{under}-estimates $e^{0.2}$, and justify
your answer using the sign of $f''$. Then check numerically.
\end{exercise}

\begin{exercise}[A guaranteed error bar]
Let $f(x) = \sqrt{x}$ and anchor at $a=9$. Use the linearization to estimate
$\sqrt{9.3}$, then use the bound $\abs{E}\le \tfrac{M}{2}h^2$ (with $M$ an upper
bound for $\abs{f''}$ on $[9,9.3]$) to produce a guaranteed interval that
contains the true value of $\sqrt{9.3}$. Verify the true value lies inside.
\end{exercise}

\begin{exercise}[Relative error in a power law]
The kinetic energy of a moving mass is $E = \tfrac12 m v^2$. If the speed $v$ is
known to within $1.5\%$ and the mass $m$ to within $0.4\%$, bound the relative
error in $E$. (Hint: take logarithms first, then differentiate.)
\end{exercise}

\begin{exercise}[Newton by hand]
Apply Newton's method to $f(x) = x^3 - 7$ to estimate $\sqrt[3]{7}$.
\begin{enumerate}
\item Write the iteration $x_{n+1} = x_n - f(x_n)/f'(x_n)$ explicitly for this
$f$.
\item Starting from $x_0 = 2$, compute $x_1$ and $x_2$ to five decimal places.
\item How many correct digits did each step gain? Relate this to the
$h\mapsto h^2$ error law.
\end{enumerate}
\end{exercise}

\begin{exercise}[When Newton stumbles]
Consider $f(x) = x^3 - 2x + 2$ with starting guess $x_0 = 0$.
\begin{enumerate}
\item Compute $x_1$ and $x_2$.
\item Show that the iteration falls into a $2$-cycle ($0 \to 1 \to 0 \to \cdots$)
and never approaches the real root near $x\approx -1.77$. Explain geometrically,
in terms of where the tangent lines point, why the method fails to converge here.
\end{enumerate}
\end{exercise}

\begin{exercise}[Conceptual: what is $\dd y$?]
A student writes: ``$\dd y$ is the change in $y$ when $x$ changes by $\dd x$.''
Correct this sentence so that it is precise. In your corrected version, explain
the difference between $\dd y$ and $\Delta y$, draw the small picture that
distinguishes them, and identify which one the derivative $\dv{y}{x}$ is a ratio
of.
\end{exercise}

\begin{exercise}[Prove the first-order property]
Let $f$ be differentiable at $a$ and let $L(x) = f(a)+f'(a)(x-a)$ be its
linearization. Prove directly from the definition of the derivative that
\[
\lim_{x\to a} \frac{f(x)-L(x)}{x-a} = 0,
\]
and explain why no line of a \emph{different} slope can satisfy this. (This is the
precise statement that the tangent line is the unique best linear approximation.)
\end{exercise}

\begin{exercise}[Proof: error formula on a polynomial]
Let $f(x) = x^3$ and anchor at $a=1$.
\begin{enumerate}
\item Write the exact error $E(x) = f(x) - L(x)$ as a polynomial in $h=x-1$ by
expanding $(1+h)^3$.
\item Show directly that $E(x) = \tfrac{f''(c)}{2}(x-1)^2$ holds, and solve for
the point $c$ as a function of $h$. Confirm that $c$ lies between $1$ and $x$, as
Theorem~\ref{ch14:thm-error} promises.
\end{enumerate}
\end{exercise}

\begin{exercise}[Linearization of a product]
Suppose $u$ and $v$ are differentiable at $a$, with values $u(a), v(a)$ and
derivatives $u'(a), v'(a)$. Using $\dd(uv) = u\,\dd v + v\,\dd u$, show that the
relative error of a product is the sum of the relative errors:
\[
\frac{\dd(uv)}{uv} = \frac{\dd u}{u} + \frac{\dd v}{v}.
\]
Interpret this as the differential of the identity $\ln(uv) = \ln u + \ln v$, and
state the analogous rule for a quotient $u/v$.
\end{exercise}
